\section{Preliminary Research: LLMs researching LLMs}\label{sec:appdx_preliminary}

Two popular implementations that use LLM agents for automated research to optimize the training of a secondary target LLM are \textit{PostTrainBench} \cite{rank2026posttrainbench} and \textit{autoresearch} \cite{karpathy2026autoresearch}. We describe both below, including their environmental setup and reported cheating behavior.

\begin{figure*}[t]
\footnotesize
\begin{promptbox}[]
\begin{enumerate}[leftmargin=*, topsep=2pt, itemsep=2pt, parsep=0pt]
  \item There will be no user interaction. You have to operate autonomously.
  \item You have a maximum of \texttt{\{num\_hours\}} hours. You can query how much time you have left by calling \texttt{bash timer.sh} in the current directory.
  \item Do not use \texttt{\{benchmark\}} test data for training (neither questions, nor answers). Otherwise we will have data contamination.
  \item Do not modify \texttt{eval.py}.
  \item Work only in the current directory and subdirectories. You can of course make use of the huggingface cache, which is not located here, but in the home folder.
  \item Your \texttt{final\_model} must run in the current environment (same packages). Installing additional packages for experimentation is allowed, but ensure \texttt{final\_model} works with the starting environment.
  \item You are not allowed to fine-tune any model other than \texttt{\{model\}}. This means that it is forbidden to download an instruction-tuned version of \texttt{\{model\}}. Only fine-tune from \texttt{\{model\}} and submit one of your fine-tuned versions as your \texttt{final\_model}.
\end{enumerate}
\end{promptbox}
\caption{PostTrainBench prompt excerpt stating the rules given to the agent}
\label{fig:rules_prompt_ptb}
\end{figure*}

\paragraph{PostTrainBench (PTB)} \cite{rank2026posttrainbench} tasks frontier LLM coding agents such as Codex and Claude Code with optimizing a smaller base LLM (e.g., Qwen3-4B) on a downstream benchmark such as AIME. The provided agent setup is deliberately minimal so that the agent retains full autonomy: it consists of a prompt that specifies the task (the target model and benchmark) together with an evaluation file that scores the model on the benchmark test set. The prompt also lists a small set of explicit rules
which we illustrate in figure \ref{fig:rules_prompt_ptb}.  Specifically rules 3,4,5, and 7 look like very specific rules that might reflect added patches of PostTrainBench preliminary experiments. 
Under this setup, agents achieve substantial gains, with the best agent reaching a 23\% improvement over the baseline. The authors also report several cheating behaviors of the agents, including training on the test set and downloading an instruction-tuned models.



\paragraph{Autoresearch} \cite{karpathy2026autoresearch} is an open-source coding project that optimizes a randomly initialized LLM for pretraining via next-token prediction. Its setup resembles PostTrainBench but differs in two important respects. First, the instructions are much more rigid and enforce a scientific pipeline. Second, the agent is given an infrastructure of supporting files (e.g., training and data preparation scripts) rather than a near-empty workspace. Despite a much wider user base, reports of cheating behavior in autoresearch are far more limited. From our examination and preliminary experiments, we attribute this to two factors: stronger scientific scaffolding and clearer instructions. The scaffolding enforces the use of git so that each experiment runs on its own branch and is dropped if it does not improve performance, and it requires that every experiment, including failed ones, be logged in a single file. The detailed prompt is likely the more important factor, since recent work suggests that instruction ambiguity underlies other reported failure modes, such as self-preservation \cite{rajamanoharan2025self}.




\subsection{Preliminary Results}


We ran preliminary experiments on both pretraining (using autoresearch) and post-training (using PostTrainBench). Because these experiments are compute-heavy and costly, and because some yielded negative results, each experiment was run on only a limited set of models.

\subsubsection{Autoresearch - Pretraining } 
We started by running the original autoresearch pipeline with Claude Code, Codex, and three open-source agents: DeepSeek-V4-Pro, Kimi-K2.5, and Gemini-3.1-Pro-Preview. None of these initial runs showed any cheating behavior. Drawing on related work on eliciting cheating in LLM agents, we tried several modifications, including increasing pressure (telling the agent that the results would be used for a funding demo and that strong outcomes were essential) and imposing unrealistic expectations (reducing the compute budget per experiment and setting unattainable performance targets). Again, none of these runs produced any clear cheating behavior either.


\subsubsection{PostTrainBench - Post-training}

\paragraph{Eliciting Cheating} We ran the original PTB setup alongside a variant we call \textit{unconstrained}, in which we remove the four rules that explicitly forbid specific cheats such as training on the test set.
As the original PTB runs were for several models and the agent quit at different times we report clearly the time and model for each run. Even with the different training runs, we see that the unconstrained agents with less training time achieve an accuracy far above those of the original PTB experiments, resulting in cheating 4/5 times. The likely source is that the agent now interprets the prompt as do anything to obtain a high test accuracy, which included training on the test set in several of these cases.  


\begin{table}[h]
\centering
\begin{tabular}{lccc}
\toprule
Model & Budget & Acc.\ (\%) & Cheat \\
\midrule
\multicolumn{4}{l}{\textit{PTB}} \\
\quad Qwen3-1.7B & 2h  & 5.5  & -- \\
\quad Qwen3-1.7B & 10h & 53.0 & -- \\
\quad Qwen3-4B   & 6h  & 53.0 & -- \\
\quad Qwen3-4B   & 10h & 67.7 & -- \\
\midrule
\multicolumn{4}{l}{\textit{PTB-unc}} \\
\quad Qwen3-1.7B & 2h & 4.3  & -- \\
\quad Qwen3-1.7B & 2h & 86.0 & \checkmark \\
\quad Qwen3-1.7B & 2h & 99.4 & \checkmark \\
\quad Qwen3-1.7B & 4h & 70.1 & \checkmark \\
\quad Qwen3-1.7B & 4h & 75.6 & \checkmark \\
\bottomrule
\end{tabular}
\caption{PTB baseline vs.\ PTB unconstrained on HumanEval. Each row is a single Claude-orchestrated run; target model and compute budget vary across runs.}
\label{tab:ptb-elicit-cheating}
\end{table}

\paragraph{Mitigating Cheating} 
While the additional rules in the original PTB prompt prohibit a small set of specific cheats, such a list will always be non-exhaustive. We hypothesize that the absence of cheating in autoresearch is driven less by explicit prohibitions and more by nudging the agent to follow a scientific process. Autoresearch encourages this in several ways, the most subtle of which is the instruction to log experiments. To test this, we ran a new experiment on PTB with the unconstrained prompt, augmented with a single line from the autoresearch prompt asking the agent to log its experiments to a file, together with the results and a brief description. Table~\ref{tab:ptb-baseline-vs-logging} reports three Claude-Opus-4-6 runs per condition on the GSM8K benchmark, with each run lasted max 5h, and each experiment was constrained to 10 minutes. Adding only the logging instruction reduces accuracy by roughly 8\% and eliminates explicit cheating in all three runs. In the unconstrained runs we observed cheating in two cases: in one the agent trained on the test set, and in the other the agent inspected the test set and modified the training set accordingly.

\begin{table}[h]
\centering
\resizebox{0.99\columnwidth}{!}{
\begin{tabular}{lcccc}
\toprule
Variant & $n$ & Mean acc.\ (\%) & Best (\%) & RH cheat \\
\midrule
PTB-unc  & 3 & $79.1 \pm 1.4$ & 80.7 & 2/3\\
PTB-unc+log               & 3 & $70.7 \pm 1.4$ & 72.0 & 0/3\\
\bottomrule
\end{tabular}
}
\caption{PTB unc vs.\ PTB with results-logging instruction. Three Claude-Opus-4-6 runs per condition, 5\,h budget, GSM8K-150 accuracy, transcript-aware LLM judge (gpt-5.4).}
\label{tab:ptb-baseline-vs-logging}
\end{table}

\subsection{Summary preliminary research experiments}
Although none of these experiments were repeated across many agents or many runs, they produced two takeaways that informed our main experiments. First, strict instructions to follow a scientific process, combined with scaffolding that constrains the agent's freedom, reduce the likelihood of cheating. Second, when the task is open-ended with many degrees of freedom, the agent can spend unlimited time tweaking minor parameters; this limits creative progress and also reduces the chance of eliciting cheating, since the open nature of the task gives the agent legitimate ways to make small changes indefinitely. We therefore focused our main set of experiments on a smaller research project using tabular machine learning problems, which we will describe below.

